\subsection{Limitations}
\label{sec:res-limitations}

Our data place several limits on what we can conclude. The Spanish records give average grades for graduates in each university, field and year. They cannot show how individual students change position, and they leave out students who do not graduate. Around half of the possible cells are unpublished, with smaller fields under-represented. We therefore simulated selection shares from published means and quartiles rather than observing actual selections. The checks in \S\ref{sec:res-selection} help assess the normal approximation, but cannot replace individual records. Admission-grade controls are also limited: they cover public universities and five broad branches, so they cannot separate the intake of fields within the same branch. Similar coefficients after adjustment therefore do not rule out differences in who enters each degree.

The Portuguese data let us follow individual students, but only at one polytechnic and into the next semester, among those with grades in both semesters. The US data cover more institutions, but the comparison has only ten areas and different grade types and years. Higher grades may reflect stronger learning as well as more generous grading. Making access more equal across degrees does not, by itself, show that we have removed grading bias or measured students' ability more accurately.
